\chapter{Classification of minimising cones}\label{ch:cones}
\module{NoCompromise.Cones.ThreeDim,
        NoCompromise.Cones.TwoDim,
        NoCompromise.Tests.Constants}

\section{The two-dimensional case}

\begin{lemma}[Structure of a minimising cone in the plane]\label{lem:cone-2d-link}
\uses{def:omega-minimal}
Let $C\subset\R^2$ be a nontrivial locally perimeter-minimising cone.  Then
the link $\Gamma:=\partial C\cap S^1$ is a finite set, $\partial C$ in the
punctured plane is exactly the union of the corresponding rays, and
$\#\Gamma$ is even.
\end{lemma}

\begin{proof}\uses{lem:bilip-chain,lem:bv-slices,def:density}
Replace $C$ by its density-one representative.  Cover the annulus
$B_2\setminus\overline{B_{1/2}}$ by finitely many polar-coordinate sectors on
which $(r,\theta)\mapsto r(\cos\theta,\sin\theta)$ and its inverse are
bi-Lipschitz.  Lemma~\ref{lem:bilip-chain} puts the pullback of $\ind C$ in
$BV$ on every such rectangle.  Cone invariance makes this pullback independent
of $r$ a.e.; Lemma~\ref{lem:bv-slices} therefore shows that its angular factor
is a $\{0,1\}$-valued $BV$ function on every coordinate arc.  Compatibility on
overlaps gives a $BV$ function on the compact circle $S^1$.

A one-dimensional $\{0,1\}$-valued $BV$ function has only finitely many jumps:
each jump has variation one, while its total variation is finite.  Its jump
set is precisely $\Gamma=\partial C\cap S^1$.  Indeed, away from a jump the
angular factor is a.e.\ constant on an arc, and cone invariance makes the
corresponding open sector one phase; at a jump both phases occupy sectors of
positive density.  Dilation then shows that the punctured boundary is exactly
the union of the rays through these jump points.  Traversing the circle, the
two values alternate at successive jumps and return to the starting value, so
the number of jumps is even.  No rectifiable-set coarea theorem or curvature
argument is used.
\end{proof}

\begin{lemma}[Two-dimensional minimising cones are halfplanes]\label{lem:cone-2d}
\uses{lem:cone-2d-link}
A nontrivial locally perimeter-minimising cone in $\R^2$ is a halfplane.
\end{lemma}

\begin{proof}\uses{lem:cone-2d-link,def:omega-minimal,lem:locality}
\emph{Two rays.}  If there are exactly two rays with angle
$\alpha\ne\pi$, replace the two radial segments in $B_1$ by the chord joining
their endpoints.  The competitor has the same trace on $\partial B_1$ and its
added interface has length $2\sin(\alpha/2)<2$, contradicting minimality.

\emph{Four or more rays.}  If there are $2m\ge4$ rays, enumerate their
endpoints cyclically and pair $1$ with $2$, $3$ with $4$, and so on.  The
corresponding chords are disjoint, each lying in the closed sector between
consecutive rays.  Fill the alternating caps so that the trace on
$\partial B_1$ is unchanged.  No chord is longer than the two radial segments
it replaces, and at least one is strictly shorter: equality would require a
paired arc of angle $\pi$, and not all paired arcs can have angle $\pi$ when
there are at least four rays.  Hence the total perimeter strictly decreases, a
contradiction.

In both comparisons the competitor is understood to agree with $C$ outside
$B_1$.  Its symmetric difference is compactly contained in, say, $B_2$;
apply local minimality there and use locality to cancel the unchanged annulus.
Thus exactly two antipodal rays occur.
\end{proof}

\section{The three-dimensional case}

\begin{lemma}[Tangent cones of a cone are cylinders]\label{lem:cone-tangent-cylinder}
\uses{prop:tangent-is-cone,lem:tangent-cone-minimizing}
Let $C\subset\R^3$ be a nontrivial locally perimeter-minimising cone and
$p\in\partial C\setminus\{0\}$.  Then every tangent limit $T$ of $C$ at $p$ is
invariant under translations in the direction $p$, and there is a measurable
$L\subset p^\perp$ with
\begin{equation}\label{eq:cone-product}
  T=L\times\R p\qquad\text{modulo null sets.}
\end{equation}
Moreover $L$ is a cone in $p^\perp$ centred at the origin.
\end{lemma}

\begin{proof}
\uses{prop:tangent-is-cone,lem:tangent-cone-minimizing,thm:monotonicity}
Blow up at $p$ along $r_j\downarrow0$.  Cone invariance gives, whenever
$1+r_jt>0$,
\begin{equation}\label{eq:cone-rescale}
  \ind C\bigl(p+r_j(z+tp)\bigr)
  =\ind C\Bigl(p+\frac{r_j}{1+r_jt}z\Bigr),
\end{equation}
and $(1+r_jt)^{-1}\to1$.  Translation and dilation continuity in $\Lloc$
therefore show $T$ is invariant under translations in $\R p$.  Choosing
coordinates with $p=e_3$, invariance means
$\ind T(x',x_3+t)=\ind T(x',x_3)$ for every $t$; Fubini and the
distributional identity $D_3\ind T=0$ give \eqref{eq:cone-product}.  That $T$
is itself a cone centred at $0$ is Proposition~\ref{prop:tangent-is-cone}
applied at $p$, whose almost-monotonicity error vanishes because $C$ is
exactly minimising; hence $L$ is a cone in $p^\perp$.
\end{proof}

\begin{lemma}[Product competitor and descent of minimality]\label{lem:cone-descent}
\uses{lem:cone-tangent-cylinder}
Let $T=L\times\R$ be locally perimeter minimising in $\R^3$.  Then $L$ is
locally perimeter minimising in $\R^2$.
\end{lemma}

\begin{proof}\uses{lem:cone-tangent-cylinder,prop:gluing,lem:bv-slices,
                    lem:locality}
Testing $D\ind T$ with product test fields on a unit vertical interval first
shows that $L$ has locally finite perimeter; equivalently this is the product
case of the BV slicing identity.
Let $A\subset\R^2$ have locally finite perimeter and
$A\triangle L\Subset B_R^2$.  For $T_0>1$ define the explicit competitor
\begin{equation}\label{eq:product-competitor}
  F_{T_0}:=\bigl(A\times(-T_0,T_0)\bigr)
    \cup\bigl(L\times((-\infty,-T_0]\cup[T_0,\infty))\bigr),
\end{equation}
so $F_{T_0}\triangle T\Subset B_R^2\times(-T_0-1,T_0+1)$.  Put
$Q_{T_0}:=B_{R+1}^2\times(-T_0-1,T_0+1)$.  The BV product and gluing formulas
give the finite local identity
\begin{equation}\label{eq:product-perimeter}
  \Per(F_{T_0};Q_{T_0})-\Per(T;Q_{T_0})
  =2T_0\bigl(\Per_2(A;B^2_{R+1})-\Per_2(L;B^2_{R+1})\bigr)
    +2|A\triangle L| .
\end{equation}
The last term is the pair of horizontal jump interfaces at heights $\pm T_0$;
it is \emph{independent of $T_0$}.  Since the symmetric difference is
compactly contained in $Q_{T_0}$, apply local minimality in an enclosing ball
and use locality to cancel the region where the two sets agree.  This gives
the corresponding inequality on $Q_{T_0}$; division by $2T_0$ and
$T_0\to\infty$ give
$\Per_2(L;B^2_{R+1})\le\Per_2(A;B^2_{R+1})$.  Since $A=L$ on the annulus
$B^2_{R+1}\setminus B^2_R$, this is exactly
$\Per_2(L;B^2_R)\le\Per_2(A;B^2_R)$.
\end{proof}

\begin{test}[Horizontal cost]\label{test:product-cost}
\eqref{eq:product-perimeter} must have horizontal cost exactly
$2|A\triangle L|$.  Getting this factor wrong (e.g.\ $|A\triangle L|$ or
$4|A\triangle L|$) still gives a $T_0$-independent term, so the limit
argument still works and the error is invisible.  A regression test on
\eqref{eq:product-perimeter} is therefore worth writing.
\end{test}

\begin{lemma}[Smoothness of the punctured cone]\label{lem:cone-smooth}
\uses{lem:cone-2d,lem:cone-descent,thm:eps-regularity}
Let $C\subset\R^3$ be a nontrivial locally perimeter-minimising cone.  Then
$\partial C\setminus\{0\}$ is a smooth embedded surface with zero mean
curvature.
\end{lemma}

\begin{proof}
\uses{lem:cone-tangent-cylinder,lem:cone-descent,lem:cone-2d,
      thm:eps-regularity,lem:tangent-cone-minimizing,thm:campanato,
      thm:quasilinear-schauder,lem:excess-ball-cyl,
      lem:fixed-normal-excess-convergence,lem:bounded-H}
By Lemmas~\ref{lem:cone-tangent-cylinder}, \ref{lem:cone-descent} and
\ref{lem:cone-2d}, every tangent cone of $C$ at every nonzero boundary point
is a halfspace; the empty and full planar factors are excluded by the
two-sided density estimate at the tangent-cone vertex.  Let $C_j$ denote the
rescalings along a subsequence converging to such a halfspace, and let $\nu$
be its normal.  Perimeter-measure convergence from
Lemma~\ref{lem:tangent-cone-minimizing} and
Lemma~\ref{lem:fixed-normal-excess-convergence}, applied to a continuous
cutoff equal to one on a smaller cylinder, show that
$\int|\nu_{C_j}-\nu|^2\,d|D\ind{C_j}|\to0$ there.  Rescaling back, the
cylindrical excess of $C$ is therefore arbitrarily small at a sufficiently
small scale around any $p\ne0$.
Theorem~\ref{thm:eps-regularity} (with $\omega=0$) then makes
$\partial C\setminus\{0\}$ a $C^{1,1/2}$ embedded surface.  Its weak mean
curvature vanishes by \eqref{eq:bounded-H} with $\omega=0$.  The
Campanato/Schauder package of Chapter~\ref{ch:elliptic}
(Theorems~\ref{thm:campanato} and \ref{thm:quasilinear-schauder}), which is
logically prior to this chapter, makes it $C^{2,\alpha}$ with zero right side, and
the mean curvature then vanishes pointwise.  Smoothness follows by iterating
the Schauder estimate to all orders.
\end{proof}

\medskip\noindent\emph{Formalisation note.} The formalisation proves the $C^{2,\alpha}$ regularity and the pointwise zero-mean-curvature equation here, proves Theorem~\ref{thm:cone-3d} from that alone (dilation invariance forces the graph to be affine), and then obtains the smoothness asserted in this lemma a posteriori: $\partial C$ is a plane.

\begin{lemma}[The link is a single great circle]\label{lem:cone-link-great-circle}
\uses{lem:cone-smooth}
With $C$ as in Lemma~\ref{lem:cone-smooth}, the link
$\Gamma:=\partial C\cap\Sph^2$ is a single great circle.
\end{lemma}

\begin{proof}\uses{lem:cone-smooth,lem:bounded-H}
By Lemma~\ref{lem:cone-smooth}, $\Gamma$ is a smooth compact embedded
one-manifold, hence a finite disjoint union of circles.  To obtain its Euler
equation, let $Y$ be a smooth tangent vector field on $\Sph^2$, let
$\eta\in C^\infty_c((1/2,2))$, and use the ambient variation
$X(r\theta):=\eta(r)Y(\theta)$ in \eqref{eq:bounded-H} with $\omega=0$.  The
first variation of the conical surface factors as
\begin{equation}\label{eq:cone-link-variation}
  0=\int_{1/2}^2\eta(r)\,dr\int_\Gamma\dvg_\Gamma Y\,ds .
\end{equation}
Choosing $\eta$ with nonzero integral shows the geodesic-curvature vector of
each component vanishes.  If a component is parametrised by unit speed
$\gamma$, its acceleration in $\R^3$ is normal to the sphere, so
$\gamma''=-\gamma$ and hence $\gamma(t)=a\cos t+b\sin t$ for orthonormal
$a,b$: its image is a great circle.  Two distinct great circles intersect,
whereas distinct components of the embedded link are disjoint.  So $\Gamma$ is
one great circle.
\end{proof}

\medskip\noindent\emph{Formalisation note.} The formalisation proves Theorem~\ref{thm:cone-3d} directly from the $C^{2,\alpha}$ regularity of Lemma~\ref{lem:cone-smooth} and derives this lemma as a corollary (local flatness of $\partial C$ plus a local-to-global great-circle argument); the endgame of the argument above, that a nonempty family of pairwise disjoint great circles is a single great circle, is formalised as well.

\begin{theorem}[Three-dimensional minimising cones are halfspaces]
\label{thm:cone-3d}
\uses{lem:cone-link-great-circle}
Every nontrivial locally perimeter-minimising cone in $\R^3$ is a halfspace.
\end{theorem}

\begin{proof}\uses{lem:cone-link-great-circle}
$\partial C$ is the plane spanned by the great circle of
Lemma~\ref{lem:cone-link-great-circle}, so $C$ is one of the two halfspaces.
\end{proof}

